% GENERATED FILE. Edit canonical lessons, then run npm run generate:books.
% canonical-source-hash: fnv1a64:e5255db91149507e
% canonical-lessons: RU-C335-ostyvat, RU-C335-razmorazhivat, RU-C335-okhlazhdat, RU-C335-rastaplivat, RU-C335-perevorachivat

\chapter{Cooling and thawing}
\label{ch:ru-cooling-and-thawing}
\hlchaptermodality{335}

\begin{chapteropening}
\textbf{By the end of this chapter:} \emph{I can describe food cooling, defrosting, chilling and melting, and turning it over.}

Complete the last lesson of chapter 335: I can describe food cooling, defrosting, chilling and melting, and turning it over.
\end{chapteropening}

\section[\texorpdfstring{\ru{остывать} (ostyvát)}{остывать (ostyvát)}]{\ru{остывать} (ostyvát) — to cool down, to get cold}

\label{lesson:RU-C335-ostyvat}



\begin{quote}
Before the new one: say the Russian for a colander, then the Russian for a sieve.
\end{quote}

\subsection*{You'll want to know: \ru{остывать}}
\textbf{\ru{остывать}} (\emph{ostyvát}) — \textquotedblleft{}to cool down, to get cold\textquotedblright{}.

\begin{quote}
\ru{Чай} \ru{уже} \ru{остыл}.
\end{quote}

\begin{quote}
\emph{Chay uzhé ostýl.}
\end{quote}

\begin{quote}
The tea has already gone cold.
\end{quote}

Imperfective; perfective \emph{\ru{остыть}} (\emph{ostýt}), \emph{\ru{остынет}}.

\begin{cousinweb}
\emph{\ru{остывать}} is built on \emph{\ru{о}-} and \emph{\ru{стыть}} \textquotedblleft{}to grow cold\textquotedblright{}, the Slavic root of \emph{\ru{студ}} \textquotedblleft{}cold\textquotedblright{} and \emph{\ru{простуда}} \textquotedblleft{}a cold\textquotedblright{}; EtymDB lists no parent.
\end{cousinweb}

\begin{grammarlens}[title={What you've built}]
One more word for telling, explaining and describing at B1.
\end{grammarlens}

\subsection*{Your turn}
\begin{itemize}
\raggedright
  \item \emph{Say it:} \emph{ostyvát}
  \item \emph{Say it:} \emph{ostyvát}, once more
  \item \emph{Say it:} a sentence of your own with \emph{ostyvát}
  \item \emph{Recall:} say the Russian for a colander, then the Russian for a sieve, then say \emph{ostyvát} again
\end{itemize}

\begin{tcolorbox}[breakable,colback=teal!4,colframe=teal!35!black,title={Before you move on}]
What does \ru{остывать} mean? (\textquotedblleft{}To cool down, to get cold\textquotedblright{}.) Say it once more.
\end{tcolorbox}
\section[\texorpdfstring{\ru{размораживать} (razmorázhivat)}{размораживать (razmorázhivat)}]{\ru{размораживать} (razmorázhivat) — to defrost, to thaw}

\label{lesson:RU-C335-razmorazhivat}



\begin{quote}
Before the new one: say the Russian for a sieve, then the Russian for to cool down.
\end{quote}

\subsection*{You'll want to know: \ru{размораживать}}
\textbf{\ru{размораживать}} (\emph{razmorázhivat}) — \textquotedblleft{}to defrost, to thaw\textquotedblright{}.

\begin{quote}
\ru{Я} \ru{разморозил} \ru{курицу}.
\end{quote}

\begin{quote}
\emph{Ya razmorózil kúritsu.}
\end{quote}

\begin{quote}
I defrosted the chicken.
\end{quote}

Imperfective; perfective \emph{\ru{разморозить}} (\emph{razmorózit}), \emph{\ru{разморожу}}.

\begin{cousinweb}
\emph{\ru{размораживать}} is built on \emph{\ru{раз}-} \textquotedblleft{}un-\textquotedblright{} and \emph{\ru{мороз}} \textquotedblleft{}frost\textquotedblright{}, inherited through Old East Slavic \emph{morozŭ} (EtymDB).
\end{cousinweb}

\begin{grammarlens}[title={What you've built}]
One more word for telling, explaining and describing at B1.
\end{grammarlens}

\subsection*{Your turn}
\begin{itemize}
\raggedright
  \item \emph{Say it:} \emph{razmorázhivat}
  \item \emph{Say it:} \emph{razmorázhivat}, once more
  \item \emph{Say it:} a sentence of your own with \emph{razmorázhivat}
  \item \emph{Recall:} say the Russian for a sieve, then the Russian for to cool down, then say \emph{razmorázhivat} again
\end{itemize}

\begin{tcolorbox}[breakable,colback=teal!4,colframe=teal!35!black,title={Before you move on}]
What does \ru{размораживать} mean? (\textquotedblleft{}To defrost, to thaw\textquotedblright{}.) Say it once more.
\end{tcolorbox}
\section[\texorpdfstring{\ru{охлаждать} (okhlazhdát)}{охлаждать (okhlazhdát)}]{\ru{охлаждать} (okhlazhdát) — to chill, to cool}

\label{lesson:RU-C335-okhlazhdat}



\begin{quote}
Before the new one: say the Russian for to cool down, then the Russian for to defrost.
\end{quote}

\subsection*{You'll want to know: \ru{охлаждать}}
\textbf{\ru{охлаждать}} (\emph{okhlazhdát}) — \textquotedblleft{}to chill, to cool\textquotedblright{}.

\begin{quote}
\ru{Сок} \ru{нужно} \ru{охладить}.
\end{quote}

\begin{quote}
\emph{Sok núzhno okhladít.}
\end{quote}

\begin{quote}
The juice needs chilling.
\end{quote}

Imperfective; perfective \emph{\ru{охладить}} (\emph{okhladít}), \emph{\ru{охлажу}}.

\begin{cousinweb}
\emph{\ru{охлаждать}} is built on \emph{\ru{о}-} and \emph{\ru{хлад}} \textquotedblleft{}cold\textquotedblright{}, the Church Slavonic shape of native \emph{\ru{холод}}, from Proto-Slavic \emph{xoldŭ} \textquotedblleft{}cold\textquotedblright{} (EtymDB).
\end{cousinweb}

\begin{grammarlens}[title={What you've built}]
One more word for telling, explaining and describing at B1.
\end{grammarlens}

\subsection*{Your turn}
\begin{itemize}
\raggedright
  \item \emph{Say it:} \emph{okhlazhdát}
  \item \emph{Say it:} \emph{okhlazhdát}, once more
  \item \emph{Say it:} a sentence of your own with \emph{okhlazhdát}
  \item \emph{Recall:} say the Russian for to cool down, then the Russian for to defrost, then say \emph{okhlazhdát} again
\end{itemize}

\begin{tcolorbox}[breakable,colback=teal!4,colframe=teal!35!black,title={Before you move on}]
What does \ru{охлаждать} mean? (\textquotedblleft{}To chill, to cool\textquotedblright{}.) Say it once more.
\end{tcolorbox}
\section[\texorpdfstring{\ru{растапливать} (rastáplivat)}{растапливать (rastáplivat)}]{\ru{растапливать} (rastáplivat) — to melt (butter); to light (a stove)}

\label{lesson:RU-C335-rastaplivat}



\begin{quote}
Before the new one: say the Russian for to defrost, then the Russian for to chill.
\end{quote}

\subsection*{You'll want to know: \ru{растапливать}}
\textbf{\ru{растапливать}} (\emph{rastáplivat}) — \textquotedblleft{}to melt (butter); to light (a stove)\textquotedblright{}.

\begin{quote}
\ru{Она} \ru{растопила} \ru{масло} \ru{в} \ru{кастрюле}.
\end{quote}

\begin{quote}
\emph{Oná rastopíla máslo v kastryúle.}
\end{quote}

\begin{quote}
She melted the butter in a saucepan.
\end{quote}

Imperfective; perfective \emph{\ru{растопить}} (\emph{rastopít}). Also \emph{\ru{растопить} \ru{печь}} — to light a stove.

\begin{cousinweb}
\emph{\ru{растапливать}} is built on \emph{\ru{рас}-} and \emph{\ru{топить}} \textquotedblleft{}to heat, to melt\textquotedblright{}, inherited from Proto-Slavic \emph{topiti} \textquotedblleft{}to heat\textquotedblright{} (EtymDB), the root of \emph{\ru{тёплый}} \textquotedblleft{}warm\textquotedblright{}.
\end{cousinweb}

\begin{grammarlens}[title={What you've built}]
One more word for telling, explaining and describing at B1.
\end{grammarlens}

\subsection*{Your turn}
\begin{itemize}
\raggedright
  \item \emph{Say it:} \emph{rastáplivat}
  \item \emph{Say it:} \emph{rastáplivat}, once more
  \item \emph{Say it:} a sentence of your own with \emph{rastáplivat}
  \item \emph{Recall:} say the Russian for to defrost, then the Russian for to chill, then say \emph{rastáplivat} again
\end{itemize}

\begin{tcolorbox}[breakable,colback=teal!4,colframe=teal!35!black,title={Before you move on}]
What does \ru{растапливать} mean? (\textquotedblleft{}To melt (butter); to light (a stove)\textquotedblright{}.) Say it once more.
\end{tcolorbox}
\section[\texorpdfstring{\ru{переворачивать} (perevoráchivat)}{переворачивать (perevoráchivat)}]{\ru{переворачивать} (perevoráchivat) — to turn over, to flip}

\label{lesson:RU-C335-perevorachivat}



\begin{quote}
Before the new one: say the Russian for to chill, then the Russian for to melt.
\end{quote}

\subsection*{You'll want to know: \ru{переворачивать}}
\textbf{\ru{переворачивать}} (\emph{perevoráchivat}) — \textquotedblleft{}to turn over, to flip\textquotedblright{}.

\begin{quote}
\ru{Он} \ru{переворачивает} \ru{блины}.
\end{quote}

\begin{quote}
\emph{On perevoráchivayet bliný.}
\end{quote}

\begin{quote}
He flips the pancakes.
\end{quote}

Imperfective; perfective \emph{\ru{перевернуть}} (\emph{perevernút}).

\begin{cousinweb}
\emph{\ru{переворачивать}} is the imperfective partner of \emph{\ru{перевернуть}} \textquotedblleft{}to turn over\textquotedblright{}; its stem \emph{-\ru{ворач}-} belongs with \emph{\ru{воротить}} \textquotedblleft{}to turn back\textquotedblright{}, Proto-Slavic \emph{vortiti}, and \emph{-\ru{вернуть}} with \emph{\ru{вертеть}} \textquotedblleft{}to turn\textquotedblright{}, Proto-Slavic \emph{vĭrtěti}, PIE \emph{wert-} (EtymDB); the two are usually seen as grades of one root.
\end{cousinweb}

\begin{grammarlens}[title={What you've built}]
One more word for telling, explaining and describing at B1.
\end{grammarlens}

\subsection*{Your turn}
\begin{itemize}
\raggedright
  \item \emph{Say it:} \emph{perevoráchivat}
  \item \emph{Say it:} \emph{perevoráchivat}, once more
  \item \emph{Say it:} a sentence of your own with \emph{perevoráchivat}
  \item \emph{Recall:} say the Russian for to chill, then the Russian for to melt, then say \emph{perevoráchivat} again
\end{itemize}

\begin{tcolorbox}[breakable,colback=teal!4,colframe=teal!35!black,title={Before you move on}]
What does \ru{переворачивать} mean? (\textquotedblleft{}To turn over, to flip\textquotedblright{}.) Say it once more.
\end{tcolorbox}
